\documentclass[11pt,a4paper]{article}
\input{../../shared/preamble}
\SpeedHeader{Geometry}{6}

\begin{document}

\SpeedTitleBlock{Daily Geometry Drill \#6 --- Answers \& Investigations}{Henry Koduthore}
\SpeedMeta{Solids and spatial reasoning}{3D Pythagoras, volumes and surface areas of prisms, pyramids, cones and spheres, Euler's formula, nets, cross-sections, frustums, inscribed solids}
\noindent\textit{New toolkit today: Space diagonal $\sqrt{a^2+b^2+c^2}$ $\cdot$ Pyramids and cones are $\tfrac13Ah$ $\cdot$ Scaling $k,k^2,k^3$ $\cdot$ $V-E+F=2$ $\cdot$ Frustum = cone minus cone.}

% ══════════════════════════════════════════════════════════════════════════════
\section*{Section A \quad Rapid Recognition}
% ══════════════════════════════════════════════════════════════════════════════
\begin{enumerate}

%── A1 ───────────────────────────────────────────────────────────────────────
\item A convex polyhedron has $12$ vertices and $30$ edges. How many faces does it have?

\ans{$20$}
\method{Euler: $V-E+F=2$, so $F=2-12+30=20$. (It is the icosahedron.)}
\inv{Each face of this polyhedron is a triangle. Check that $3F=2E$, and find how many faces meet at each vertex.}

%── A2 ───────────────────────────────────────────────────────────────────────
\item Find the length of a space diagonal of a $3\times4\times12$ cuboid.

\ans{$13$}
\method{Apply Pythagoras twice: $\sqrt{3^2+4^2+12^2}=\sqrt{169}=13$.}
\inv{Which whole numbers from $1$ to $10$ can be the space diagonal of a cuboid with integer edges?}

%── A3 ───────────────────────────────────────────────────────────────────────
\item Find the volume of a cone with base radius $3$ and height $4$.

\ans{$12\pi$}
\method{$\tfrac13\pi r^2h=\tfrac13\pi\cdot9\cdot4=12\pi$: one third of the cylinder with the same base and height.}
\inv{Find the volume of the solid made by spinning a $3$--$4$--$5$ triangle about its hypotenuse.}

%── A4 ───────────────────────────────────────────────────────────────────────
\item A sphere has surface area $36\pi$. Find its volume.

\ans{$36\pi$}
\method{$4\pi r^2=36\pi$ gives $r=3$, and then $\tfrac43\pi r^3=36\pi$.}
\inv{Show that $r=3$ is the only radius for which a sphere's volume and surface area are the same number. For which edge is the same true of a cube?}

%── A5 ───────────────────────────────────────────────────────────────────────
\item Find the distance between the points $(1,2,3)$ and $(3,5,9)$. (It is Day 1's distance formula with one more square.)

\ans{$7$}
\method{$\sqrt{2^2+3^2+6^2}=\sqrt{49}=7$: the same as the space diagonal of a $2\times3\times6$ box.}
\inv{Find a point $P$ on the $x$-axis equidistant from $(1,2,3)$ and $(3,5,9)$.}

%── A6 ───────────────────────────────────────────────────────────────────────
\item A cube has a space diagonal of length $6$. Find its total surface area.

\ans{$72$}
\method{The space diagonal is $a\sqrt3$, so $a^2=\tfrac{36}{3}=12$ and the surface area is $6a^2=72$. There is no need to find $a$ itself.}
\inv{What is the total length of the cube's twelve edges?}

%── A7 ───────────────────────────────────────────────────────────────────────
\item Two similar solids have corresponding lengths in the ratio $2:3$. The volume of the smaller is what fraction of the volume of the larger?

\ans{$\dfrac{8}{27}$}
\method{Lengths scale by $k=\tfrac23$, areas by $k^2$ and volumes by $k^3=\tfrac{8}{27}$.}
\inv{The same two solids are painted. What fraction of the larger one's paint does the smaller one need?}

%── A8 ───────────────────────────────────────────────────────────────────────
\item Find the curved surface area of a cone with base radius $5$ and height $12$.

\ans{$65\pi$}
\method{The slant height is $\sqrt{5^2+12^2}=13$, and the curved surface is $\pi rl=65\pi$.}
\inv{Unroll the curved surface into a sector. What is the sector's angle?}

%── A9 ───────────────────────────────────────────────────────────────────────
\item Find the total surface area of a regular tetrahedron with every edge of length $2$.

\ans{$4\sqrt3$}
\method{Four equilateral triangles, each $\tfrac{\sqrt3}{4}\cdot2^2=\sqrt3$.}
\inv{Find the height of the tetrahedron (the distance from a vertex to the opposite face).}

%── A10 ──────────────────────────────────────────────────────────────────────
\item A regular octahedron has $8$ faces, all of them triangles. How many edges does it have?

\ans{$12$}
\method{Count the face--edge pairs: $8\cdot3=24$, and each edge lies on two faces, so $E=12$. (Check: $V=2-F+E=6$.)}
\inv{The centres of the faces of a cube form an octahedron. Which solid do the centres of the faces of an octahedron form?}

\end{enumerate}

% ══════════════════════════════════════════════════════════════════════════════
\section*{Section B \quad Manipulation Drills}
% ══════════════════════════════════════════════════════════════════════════════
\begin{enumerate}

%── B1 ───────────────────────────────────────────────────────────────────────
\item A sector with angle $\theta$ is cut from a disc of radius $6$ and rolled up, without overlap, into the curved surface of a cone of height $4\sqrt2$. What is $\theta$? \textit{\small(after Primer 1.4 P28)}
\begin{itemize}
\item[A)] $60^\circ$
\item[B)] $90^\circ$
\item[C)] $120^\circ$
\item[D)] $150^\circ$
\item[E)] $180^\circ$
\end{itemize}

\ans{C) $120^\circ$}
\method{The slant height is the disc's radius $6$, so the cone's base radius is $\sqrt{36-32}=2$. The sector's arc becomes the base circumference: $\tfrac{\theta}{360^\circ}\cdot12\pi=4\pi$, so $\theta=120^\circ$.}
\inv{Which angle $\theta$ gives the cone of greatest volume? Maximise $r^2h$ subject to $r^2+h^2=36$.}

%── B2 ───────────────────────────────────────────────────────────────────────
\item The base edges of a square-based pyramid are shortened by $10\%$. By what factor must its height be multiplied to keep its volume the same? \textit{\small(after Primer 1.4 P23)}
\begin{itemize}[itemsep=3pt]
\item[A)] $\dfrac{10}{9}$
\item[B)] $\dfrac{100}{81}$
\item[C)] $\dfrac{11}{10}$
\item[D)] $\dfrac{121}{100}$
\item[E)] $\dfrac{81}{100}$
\end{itemize}

\ans{B) $\dfrac{100}{81}$}
\method{The volume is $\tfrac13s^2h$. The base area is multiplied by $0.9^2=\tfrac{81}{100}$, so the height must be multiplied by $\tfrac{100}{81}$ (an increase of about $23\%$). Option A undoes the change to one length only.}
\inv{If the height rises by $p\%$ while the base edges fall by $p\%$, show that the volume always falls when $0<p<100$.}

%── B3 ───────────────────────────────────────────────────────────────────────
\item A cuboid is built from unit cubes, painted on the outside, then taken apart. Exactly $12$ of the unit cubes have no paint on them. What is the smallest possible number of unit cubes in the cuboid? \textit{\small(after Primer 8.8 P12)}
\begin{itemize}
\item[A)] $80$
\item[B)] $90$
\item[C)] $96$
\item[D)] $126$
\item[E)] $60$
\end{itemize}

\ans{A) $80$}
\method{For an $a\times b\times c$ block the unpainted cubes form an $(a-2)\times(b-2)\times(c-2)$ block, so $(a-2)(b-2)(c-2)=12$. The factorisations $1\cdot1\cdot12$, $1\cdot2\cdot6$, $1\cdot3\cdot4$ and $2\cdot2\cdot3$ give $3\cdot3\cdot14=126$, $3\cdot4\cdot8=96$, $3\cdot5\cdot6=90$ and $4\cdot4\cdot5=80$. The most balanced block is the smallest.}
\inv{For the $4\times4\times5$ block, how many unit cubes have exactly one painted face, and how many have exactly two?}

%── B4 ───────────────────────────────────────────────────────────────────────
\item A pyramid has an equilateral-triangle base of side $6$, and its three sloping edges all have length $b$. Its volume is $18$. Find $b$. \textit{\small(after Hercules Set 2 Paper 2 Q13)}
\begin{itemize}
\item[A)] $6$
\item[B)] $4\sqrt2$
\item[C)] $\sqrt{15}$
\item[D)] $2\sqrt6$
\item[E)] $2\sqrt3$
\end{itemize}

\ans{D) $2\sqrt6$}
\method{The base has area $9\sqrt3$, so $\tfrac13\cdot9\sqrt3\cdot h=18$ and $h=2\sqrt3$. Equal sloping edges put the apex above the base's centre, which is the circumradius $2\sqrt3$ from each vertex. So $b^2=12+12=24$. Option C uses the inradius $\sqrt3$ instead.}
\inv{Find the volume in terms of $a$ and $b$ for base side $a$, and check that it gives $18$ here.}

%── B5 ───────────────────────────────────────────────────────────────────────
\item A cube is cut into $n^3$ equal smaller cubes. The pieces have $5$ times the total surface area of the original cube. How many pieces are there? \textit{\small(after OxbridgeApplications TMUA Mock Exam 2 Q30)}

\ans{$125$}
\method{Each piece has edge $\tfrac an$, so the total area is $n^3\cdot6\left(\tfrac an\right)^2=6a^2\cdot n$. The factor is $n=5$, giving $125$ pieces.}
\inv{A $2\times3\times4$ cuboid is cut into unit cubes. By what factor does the total surface area grow?}

%── B6 ───────────────────────────────────────────────────────────────────────
\item \begin{minipage}[t]{0.60\linewidth}
The net shown is folded into a die, and opposite faces of the die add up to $7$. Three faces are already numbered. Which number belongs on the shaded face?
\end{minipage}\hfill
\begin{minipage}[t]{0.36\linewidth}
\vspace{0pt}
\centering
\begin{tikzpicture}[scale=0.55]
\fill[gray!30] (2,-2) rectangle (3,-1);
\foreach \x/\y in {0/0,1/0,1/-1,2/-1,2/-2,3/-2} \draw[speedgeom] (\x,\y) rectangle ++(1,1);
\node at (0.5,0.5) {$1$}; \node at (1.5,0.5) {$2$}; \node at (1.5,-0.5) {$3$};
\end{tikzpicture}
\end{minipage}

\ans{$5$}
\method{Fold the square marked $3$ down as the base. The $2$ stands up behind it, the square to the right of $3$ becomes the right wall, the shaded square, hinged to that wall's front edge, swings round to become the front, the last square becomes the top and the $1$ becomes the left wall. So the shaded face is opposite the $2$ and carries $7-2=5$.}
\inv{Of the eleven nets of a cube, how many contain a straight row of four squares?}

%── B7 ───────────────────────────────────────────────────────────────────────
\item A sphere is inscribed in a cube, and a smaller cube is inscribed in that sphere. What is the ratio of the larger cube's volume to the smaller cube's? \textit{\small(after Beyond Horizon Set 4 Paper 1 Q2)}
\begin{itemize}
\item[A)] $3$
\item[B)] $2\sqrt2$
\item[C)] $9$
\item[D)] $27$
\item[E)] $3\sqrt3$
\end{itemize}

\ans{E) $3\sqrt3$}
\method{If the large cube has edge $a$, the sphere has diameter $a$. That diameter is the small cube's space diagonal, so the small cube has edge $\tfrac{a}{\sqrt3}$. The ratio of volumes is $(\sqrt3)^3=3\sqrt3$.}
\inv{Do the same with a square, a circle and a square: the ratio of the areas is $2$. What is it in $n$ dimensions?}

%── B8 ───────────────────────────────────────────────────────────────────────
\item Regular pentagons meet three at every vertex of a convex polyhedron. The \emph{angle deficit} at a vertex is $360^\circ$ minus the face angles there, and for any convex polyhedron the deficits add up to $720^\circ$. How many faces does this polyhedron have?

\ans{$12$}
\method{Each angle of a regular pentagon is $108^\circ$ (Day 4), so the deficit is $360^\circ-3\cdot108^\circ=36^\circ$ and $V=\tfrac{720}{36}=20$. Each face has $5$ vertices and each vertex is on $3$ faces, so $F=\tfrac{3\cdot20}{5}=12$. Check: $E=30$ and $20-30+12=2$.}
\inv{Repeat for squares three at a vertex, and for equilateral triangles three, four or five at a vertex. Why can regular hexagons never meet three at a vertex of a convex polyhedron?}

%── B9 ───────────────────────────────────────────────────────────────────────
\item A solid cone and a solid hemisphere have the same base radius $r$ and the same volume. What is the height of the cone?
\begin{itemize}[itemsep=3pt]
\item[A)] $2r$
\item[B)] $r$
\item[C)] $\dfrac32r$
\item[D)] $4r$
\item[E)] $r\sqrt2$
\end{itemize}

\ans{A) $2r$}
\method{$\tfrac13\pi r^2h=\tfrac23\pi r^3$ gives $h=2r$. Option D matches the cone with a whole sphere.}
\inv{Glue the cone to the hemisphere, base to base, to make a spinning top. Find its total surface area.}

%── B10 ──────────────────────────────────────────────────────────────────────
\item Each corner of a cube is sliced off by a small plane cut, so that each corner is replaced by a small triangle. How many edges does the new solid have?
\begin{itemize}
\item[A)] $24$
\item[B)] $30$
\item[C)] $48$
\item[D)] $36$
\item[E)] $42$
\end{itemize}

\ans{D) $36$}
\method{There are $6+8=14$ faces. Each corner becomes $3$ vertices, so $V=24$, and $3$ edges meet at every vertex, so $E=\tfrac{3\cdot24}{2}=36$. Check: $24-36+14=2$.}
\inv{Now make the cuts so deep that neighbouring triangles touch at the midpoints of the edges. Find $V$, $E$ and $F$ for this solid (the cuboctahedron).}

\end{enumerate}

% ══════════════════════════════════════════════════════════════════════════════
\section*{Section C \quad Substitution \& Structure}
% ══════════════════════════════════════════════════════════════════════════════
\begin{enumerate}

%── C1 ───────────────────────────────────────────────────────────────────────
\item A cuboid has integer edge lengths, and all eight of its vertices lie on a sphere of radius $3$. What is its volume? \textit{\small(after SMC 2016 Q23)}
\begin{itemize}
\item[A)] $36$
\item[B)] $32$
\item[C)] $24$
\item[D)] $27$
\item[E)] $48$
\end{itemize}

\ans{B) $32$}
\method{A space diagonal is a diameter, so $a^2+b^2+c^2=36$. Using the squares $1,4,9,16,25$, the only way to make $36$ is $4+16+16$. So the edges are $2,4,4$ and the volume is $32$.}
\inv{Among \emph{all} cuboids with space diagonal $6$, show that the cube has the largest volume, $24\sqrt3$, which beats $32$.}

%── C2 ───────────────────────────────────────────────────────────────────────
\item \begin{minipage}[t]{0.60\linewidth}
A cube of edge $2$ has top corner $T$ and, below it, bottom corner $B$. $M$ and $N$ are the midpoints of the two top edges at $T$, and $P$ and $Q$ are the two bottom corners next to $B$. The plane through $M$, $N$, $Q$ and $P$ cuts the cube in two. What fraction of the cube is the piece containing $T$? \textit{\small(after SMC 2023 Q16)}
\begin{itemize}[itemsep=3pt]
\item[A)] $\dfrac14$
\item[B)] $\dfrac13$
\item[C)] $\dfrac{5}{16}$
\item[D)] $\dfrac{5}{24}$
\item[E)] $\dfrac{7}{24}$
\end{itemize}
\end{minipage}\hfill
\begin{minipage}[t]{0.36\linewidth}
\vspace{0pt}
\centering
\begin{tikzpicture}[scale=1.0]
\fill[gray!25] (1,2) -- (2.45,2.3) -- (2.9,0.6) -- (0,0) -- cycle;
\draw[speedgeom] (0,0) -- (2,0) -- (2,2) -- (0,2) -- cycle;
\draw[speedgeom] (2,0) -- (2.9,0.6) -- (2.9,2.6) -- (2,2); \draw[speedgeom] (0,2) -- (0.9,2.6) -- (2.9,2.6);
\draw[speedaux] (0,0) -- (0.9,0.6) -- (2.9,0.6) (0.9,0.6) -- (0.9,2.6) (0,0) -- (2.9,0.6);
\draw[speedgeom] (0,0) -- (1,2) -- (2.45,2.3) -- (2.9,0.6);
\node[above left] at (1,2) {\small $M$}; \node[above] at (2.45,2.3) {\small $N$}; \node[below right] at (2,2) {\small $T$};
\node[below left] at (0,0) {\small $P$}; \node[right] at (2.9,0.6) {\small $Q$}; \node[below] at (2,0) {\small $B$};
\end{tikzpicture}
\end{minipage}

\ans{E) $\dfrac{7}{24}$}
\method{The piece is a frustum: its top is the triangle $MNT$ (area $\tfrac12$), its base is the triangle $PQB$ (area $2$) and its height is $2$. Extend the sloping edges up to the apex: the full pyramid on $PQB$ has height $4$ and volume $\tfrac13\cdot2\cdot4=\tfrac83$, and the missing top pyramid is half the size in every direction, so its volume is $\tfrac18\cdot\tfrac83=\tfrac13$. The piece is $\tfrac83-\tfrac13=\tfrac73$ out of $8$. Option C treats it as a prism with the average cross-section.}
\inv{Move $M$ and $N$ along the top edges, staying the same distance $t$ from $T$. For which $t$ does the plane through $M$, $N$, $Q$ and $P$ cut the cube exactly in half?}

%── C3 ───────────────────────────────────────────────────────────────────────
\item Two cubes of edge $1$ have the same centre and the same vertical axis of symmetry, and one is turned through $45^\circ$ about that axis. What is the volume of the solid the two cubes form together? \textit{\small(after SMC 2011 Q25)}
\begin{itemize}[itemsep=3pt]
\item[A)] $4-2\sqrt2$
\item[B)] $2\sqrt2-1$
\item[C)] $3-\sqrt2$
\item[D)] $2\sqrt2-2$
\item[E)] $1+\dfrac{\sqrt2}{4}$
\end{itemize}

\ans{A) $4-2\sqrt2$}
\method{Both cubes fill the same height $1$, so only the plan view matters. The turned square's four corners stick out past the sides of the other square by $\tfrac{\sqrt2}{2}-\tfrac12$, and each corner that sticks out is a right-angled isosceles triangle of area $\left(\tfrac{\sqrt2-1}{2}\right)^2$. Four of them total $3-2\sqrt2$, so the overlap is $1-(3-2\sqrt2)=2\sqrt2-2$ (a regular octagon, as in Day 4). Add the two volumes and subtract the overlap: $2-(2\sqrt2-2)=4-2\sqrt2$. Option D is the overlap alone.}
\inv{Find the total surface area of the combined solid.}

%── C4 ───────────────────────────────────────────────────────────────────────
\item A hollow cone with base radius $3$ and height $4$ stands on its vertex with its rim level. A ball inside it touches the cone all the way round, and the top of the ball is level with the rim. What fraction of the cone's volume does the ball fill? \textit{\small(after Beyond Horizon Set 2 Paper 1 Q11)}
\begin{itemize}[itemsep=3pt]
\item[A)] $\dfrac12$
\item[B)] $\dfrac{9}{32}$
\item[C)] $\dfrac38$
\item[D)] $\dfrac23$
\item[E)] $\dfrac14$
\end{itemize}

\ans{C) $\dfrac38$}
\method{Cut through the axis. The cone becomes a triangle with sides $5,5,6$, and the ball's great circle is its incircle, so the radius is $\tfrac{\text{area}}{\text{semi-perimeter}}=\tfrac{12}{8}=\tfrac32$. The ball has volume $\tfrac43\pi\cdot\tfrac{27}{8}=\tfrac{9\pi}{2}$ and the cone $12\pi$, a ratio of $\tfrac38$.}
\inv{A second, smaller ball sits in the tip of the cone, touching the first ball and the cone. Find its radius, using Day 5's circles in an angle.}

%── C5 ───────────────────────────────────────────────────────────────────────
\item A cone has base radius $3\sqrt2$ and height $6$. A cube stands on the cone's base with its four top corners on the curved surface. What is the edge of the cube? \textit{\small(after JZMaths TMUA Mock Set B Paper 1 Q18)}
\begin{itemize}[itemsep=3pt]
\item[A)] $2$
\item[B)] $2\sqrt2$
\item[C)] $12-6\sqrt2$
\item[D)] $3$
\item[E)] $\dfrac{3\sqrt2}{2}$
\end{itemize}

\ans{D) $3$}
\method{At height $s$ the cone's radius is $3\sqrt2\left(1-\tfrac s6\right)$. The top corners are half a diagonal, $\tfrac{s}{\sqrt2}$, from the axis. So $\tfrac{s}{\sqrt2}=3\sqrt2\left(1-\tfrac s6\right)$, i.e.\ $s=6-s$, and $s=3$. Option C uses half the edge instead of half the diagonal.}
\inv{Replace the cube by a cylinder standing on the base. Show that the largest possible cylinder fills $\tfrac49$ of the cone.}

%── C6 ───────────────────────────────────────────────────────────────────────
\item A regular tetrahedron and a cube are inscribed in the same sphere. What is the ratio of the tetrahedron's volume to the cube's? \textit{\small(after Primer 9.1 P41)}
\begin{itemize}[itemsep=3pt]
\item[A)] $\dfrac16$
\item[B)] $\dfrac13$
\item[C)] $\dfrac{\sqrt2}{4}$
\item[D)] $\dfrac29$
\item[E)] $\dfrac12$
\end{itemize}

\ans{B) $\dfrac13$}
\method{Four alternate corners of a cube are joined by face diagonals of equal length, so they form a regular tetrahedron with the same circumsphere as the cube. A regular tetrahedron is fixed by its circumradius, so this is the tetrahedron in the question. It is the cube minus four corner pyramids, each $\tfrac13\cdot\tfrac12a^2\cdot a=\tfrac16a^3$, leaving $a^3-\tfrac46a^3=\tfrac13a^3$.}
\inv{Add the regular octahedron inscribed in the same sphere. Show that its volume is $\tfrac{\sqrt3}{2}$ of the cube's.}

%── C7 ───────────────────────────────────────────────────────────────────────
\item A bucket is a frustum of a cone with base radius $1$, top radius $2$ and height $3$. It is filled with water to half its height. What fraction of the bucket is full? \textit{\small(after Beyond Horizon Set 1 Paper 1 Q20)}
\begin{itemize}[itemsep=3pt]
\item[A)] $\dfrac12$
\item[B)] $\dfrac38$
\item[C)] $\dfrac{19}{56}$
\item[D)] $\dfrac{27}{56}$
\item[E)] $\dfrac13$
\end{itemize}

\ans{C) $\dfrac{19}{56}$}
\method{Complete the cone. The radius grows by $1$ for every $3$ of height, so the apex is $3$ below the base. Volumes of the cones cut off at the apex are proportional to radius cubed: the bucket is $2^3-1^3=7$ units and the water $1.5^3-1^3=\tfrac{19}{8}$, so the fraction is $\tfrac{19}{56}$. Option D forgets to remove the missing tip.}
\inv{How deep must the water be to fill exactly half of the bucket?}

%── C8 ───────────────────────────────────────────────────────────────────────
\item Every face of a convex polyhedron is a triangle, and at every vertex either five or six faces meet. How many vertices have exactly five faces?
\begin{itemize}
\item[A)] $20$
\item[B)] $6$
\item[C)] $10$
\item[D)] $24$
\item[E)] $12$
\end{itemize}

\ans{E) $12$}
\method{Let there be $V_5$ and $V_6$ vertices of each kind. Triangles give $3F=2E$, and Euler gives $V=2+E-F=2+\tfrac E3$. Counting edge ends gives $5V_5+6V_6=2E$, so $V_5=6V-2E=6\left(2+\tfrac E3\right)-2E=12$, whatever the number of six-fold vertices.}
\inv{Swap faces and vertices: a polyhedron with pentagonal and hexagonal faces, three at each vertex. Show that it has exactly $12$ pentagons (a football).}

\end{enumerate}

% ══════════════════════════════════════════════════════════════════════════════
\section*{Section D \quad Challenge}
% ══════════════════════════════════════════════════════════════════════════════
\begin{enumerate}

%── D1 ───────────────────────────────────────────────────────────────────────
\item Beatrix draws on the faces of a cube of edge $2$. On each face she either draws nothing or draws one straight line joining the midpoints of two edges of that face (adjacent or opposite edges). She wants her lines to form a single closed loop. What is the greatest possible length of the loop? \textit{\small(after SMC 2014 Q18)}
\begin{itemize}
\item[A)] $4+4\sqrt2$
\item[B)] $8+2\sqrt2$
\item[C)] $6+3\sqrt2$
\item[D)] $6\sqrt2$
\item[E)] $8$
\end{itemize}

\ans{A) $4+4\sqrt2$}
\method{Where the loop crosses from one face to the next, it passes the midpoint of the edge the two faces share. So a loop on $k$ faces is a cycle of faces $F_1,\dots,F_k$ in which neighbours are adjacent faces, and the line on $F_i$ is long ($2$, joining opposite edges) exactly when $F_{i-1}$ and $F_{i+1}$ are opposite faces; otherwise it is $\sqrt2$. On all six faces, each pair of opposite faces sits $2$ or $3$ places apart round the cycle (not $1$: opposite faces do not touch), and only a pair $2$ apart makes a long line. A pair $2$ apart has positions of the same parity, and the three odd positions cannot be paired among themselves, so at most two pairs are $2$ apart: the length is at most $2\cdot2+4\sqrt2$. Top, front, bottom, left, back, right achieves it. On five faces there are exactly two long lines ($4+3\sqrt2$), and on four at most $8$. Option B is the longest \emph{open} line, which cannot be closed.}
\inv{The source asks for the longest unbroken line that need not close up, and the answer is $8+2\sqrt2$. Which faces does it use, and why can it not be closed?}

%── D2 ───────────────────────────────────────────────────────────────────────
\item A cylinder is inscribed in a sphere, with both rims on the sphere. What is the largest possible area of its curved surface, as a fraction of the sphere's surface area? \textit{\small(after nabla-tuition TMUA Paper 1 Q18)}
\begin{itemize}[itemsep=3pt]
\item[A)] $\dfrac13$
\item[B)] $\dfrac23$
\item[C)] $\dfrac{1}{\sqrt2}$
\item[D)] $\dfrac12$
\item[E)] $\dfrac34$
\end{itemize}

\ans{D) $\dfrac12$}
\method{Let the sphere have radius $R$ and the cylinder radius $r$ and height $h$, so $r^2+\left(\tfrac h2\right)^2=R^2$. By AM--GM, $r\cdot\tfrac h2\le\tfrac12\left(r^2+\tfrac{h^2}{4}\right)=\tfrac{R^2}{2}$, so the curved area is $2\pi rh=4\pi r\cdot\tfrac h2\le2\pi R^2$: half of $4\pi R^2$. Equality holds when $r=\tfrac h2=\tfrac{R}{\sqrt2}$.}
\inv{The source maximises the cylinder's \emph{volume} instead. Is it the same cylinder? Find the height for each.}

%── D3 ───────────────────────────────────────────────────────────────────────
\item A tetrahedron $ABCD$ has $AB=CD=\sqrt5$, $AC=BD=\sqrt{10}$ and $AD=BC=\sqrt{13}$. What is its volume? \textit{\small(after SMC 2022 Q25)}
\begin{itemize}
\item[A)] $1$
\item[B)] $3$
\item[C)] $2$
\item[D)] $6$
\item[E)] $4$
\end{itemize}

\ans{C) $2$}
\method{The three lengths are the face diagonals of a $1\times2\times3$ box: $1+4=5$, $1+9=10$, $4+9=13$. Four alternate corners of the box, such as $(0,0,0)$, $(1,2,0)$, $(1,0,3)$ and $(0,2,3)$, form this tetrahedron, with opposite edges on opposite faces of the box. It is the box minus four corner pyramids, each $\tfrac16\cdot1\cdot2\cdot3=1$, so the volume is $6-4=2$.}
\inv{Show that the box exists (with real edges) exactly when every face of the tetrahedron is acute, the fact D5 proves directly.}

%── D4 ───────────────────────────────────────────────────────────────────────
\item A plane perpendicular to a space diagonal of a cube of edge $2$ slides from one end of that diagonal to the other. What is the largest area of cross-section it makes? \textit{\small(after SMC 2019 Q23)}
\begin{itemize}
\item[A)] $2\sqrt3$
\item[B)] $3\sqrt3$
\item[C)] $4\sqrt2$
\item[D)] $6$
\item[E)] $3\sqrt2$
\end{itemize}

\ans{B) $3\sqrt3$}
\method{Take the cube $[0,2]^3$ and the planes $x+y+z=s$. For $0\le s\le2$ the section is an equilateral triangle of side $s\sqrt2$, with area $\tfrac{\sqrt3}{2}s^2$. For $2\le s\le4$ it is that triangle with three corners of side $(s-2)\sqrt2$ cut off, a hexagon of area $\tfrac{\sqrt3}{2}\left(s^2-3(s-2)^2\right)$. This is largest at $s=3$, the centre, where it is the regular hexagon through six edge midpoints: $\tfrac{\sqrt3}{2}(9-3)=3\sqrt3$. Option C is the rectangle through two opposite edges, which is not perpendicular to a space diagonal.}
\inv{The largest section of the cube by \emph{any} plane is that rectangle, of area $4\sqrt2$. Check that it beats the hexagon, and show that its diagonals are two of the cube's space diagonals.}

%── D5 ───────────────────────────────────────────────────────────────────────
\item \begin{minipage}[t]{0.60\linewidth}
A tetrahedron $ABCD$ has $AB=CD$, $AC=BD$ and $AD=BC$. Prove that every face of $ABCD$ is an acute-angled triangle.
\end{minipage}\hfill
\begin{minipage}[t]{0.36\linewidth}
\vspace{0pt}
\centering
\begin{tikzpicture}[scale=0.7]
\draw[speedgeom] (0,0) -- (2.4,3) -- (4,0) -- (3,-1.4) -- cycle; \draw[speedgeom] (2.4,3) -- (3,-1.4);
\draw[speedaux] (0,0) -- (4,0) (2.4,3) -- (2,0) -- (3,-1.4);
\node[left] at (0,0) {\small $A$}; \node[right] at (4,0) {\small $B$}; \node[above] at (2.4,3) {\small $C$};
\node[below] at (3,-1.4) {\small $D$}; \node[above left] at (2,0) {\small $M$};
\end{tikzpicture}
\end{minipage}

\ans{Proof: see method}
\method{Take the face angle $\angle ACB$. Let $M$ be the midpoint of $AB$. Triangles $ABC$ and $BAD$ are congruent ($AB$ is common, $AC=BD$, $BC=AD$), so their medians to $AB$ are equal: $CM=DM$. (By Apollonius, $4CM^2=2CA^2+2CB^2-AB^2=2DB^2+2DA^2-AB^2=4DM^2$.) $C$, $M$ and $D$ are not collinear, or all four points would lie in one plane, so the triangle inequality in $CMD$ is strict: $CD<CM+MD=2CM$. Since $CD=AB$, this gives $CM>\tfrac12AB=MA=MB$, so $C$ lies outside the circle with diameter $AB$ and $\angle ACB<90^\circ$. Every face angle is opposite an edge of its face, and the hypotheses are symmetric in the three pairs of opposite edges, so the same argument applies to all twelve face angles. $\blacksquare$}
\inv{Conversely, show that any acute triangle can be folded into such a tetrahedron: take the triangle's medial triangle as the base and fold the three corners up until they meet.}

\end{enumerate}

\section*{Top 5 Patterns --- Worksheet \#6}
\begin{enumerate}[leftmargin=2em, itemsep=4pt]
  \item \textbf{Pythagoras twice.} A space diagonal is $\sqrt{a^2+b^2+c^2}$, and a sphere round a box has the box's diagonal as its diameter. \seealso{A2, A5, A6, B7, C1}
  \item \textbf{One third of base times height, then scale.} Pyramids and cones are $\tfrac13Ah$; similar solids scale volume by $k^3$. \seealso{A3, A7, B2, B4, C7}
  \item \textbf{The section through the axis.} A cone, ball or cylinder becomes a triangle, circle or rectangle in the plane. \seealso{A8, B1, C4, C5, D2}
  \item \textbf{Count twice.} Face--edge and vertex--edge pairs, with $V-E+F=2$, pin down a polyhedron. \seealso{A1, A10, B8, B10, C8}
  \item \textbf{Cut and subtract.} A frustum is a cone minus a cone; a box minus four corners is a tetrahedron; a union is the sum minus the overlap. \seealso{C2, C3, C6, D3, D4}
\end{enumerate}

\section*{Common Traps --- Worksheet \#6}
\begin{enumerate}[leftmargin=2em, itemsep=4pt]
  \item \textbf{Answering the overlap or the whole.} C3 asks for the union and C7 for part of a frustum; the overlap and the full cone are one step short. \seealso{C3, C7}
  \item \textbf{Half the edge where half the diagonal is needed.} A square's corners are $\tfrac{s}{\sqrt2}$ from its centre; a triangle's vertices are the circumradius from its centre. \seealso{C5, B4}
  \item \textbf{Scaling linearly.} Areas scale by $k^2$ and volumes by $k^3$, so a $10\%$ cut in length is a $19\%$ cut in area. \seealso{B2, A7, B5}
  \item \textbf{Reusing a remembered answer after the question has changed.} The longest open line is not the longest closed loop, and the biggest section of a cube is not the biggest one perpendicular to a diagonal. \seealso{D1, D4}
  \item \textbf{Treating a tapering piece as a prism.} Average area times height overestimates a frustum; complete the cone and subtract. \seealso{C2, C7}
\end{enumerate}

\SpeedClosing{``In three dimensions, find the flat picture: the section through the axis, the net, or the triangle hidden inside the solid.''}

\end{document}
